\documentclass[10pt,a4paper]{amsart}
\usepackage{amsmath,amssymb,amsthm,graphicx}
% Retained minimal preamble for the Stacks Project chapter more-algebra.tex (More on Algebra).
% The chapter file pinned-original-source.tex carries only \input{preamble} on its first line, so no
% preamble text lives inside the chapter itself.  The Stacks Project uses one single shared preamble for
% every chapter; the snapshot stages it identically as sources/stacks-chow/preamble.tex and
% sources/stacks-etale-cohomology/preamble.tex.  The block below is copied mechanically, line for line,
% from sources/stacks-chow/preamble.tex lines 14-18 (the Xy-pic setup) and 154-217 (the native theorem
% declarations together with the mathematical macros).  Nothing was added, removed or redefined; the
% wrapper class and the amsmath/amssymb/amsthm/graphicx base packages are supplied by the assembler.
% Chapter: more-algebra.tex -- More on Algebra.
% For commutative diagrams we use Xy-pic
\usepackage[all]{xy}

% We use 2cell for 2-commutative diagrams.
\xyoption{2cell}

% Theorem environments.
%
\theoremstyle{plain}
\newtheorem{theorem}[subsection]{Theorem}
\newtheorem{proposition}[subsection]{Proposition}
\newtheorem{lemma}[subsection]{Lemma}

\theoremstyle{definition}
\newtheorem{definition}[subsection]{Definition}
\newtheorem{example}[subsection]{Example}
\newtheorem{exercise}[subsection]{Exercise}
\newtheorem{situation}[subsection]{Situation}

\theoremstyle{remark}
\newtheorem{remark}[subsection]{Remark}
\newtheorem{remarks}[subsection]{Remarks}

\numberwithin{equation}{subsection}

% Macros
%
\def\lim{\mathop{\mathrm{lim}}\nolimits}
\def\colim{\mathop{\mathrm{colim}}\nolimits}
\def\Spec{\mathop{\mathrm{Spec}}}
\def\Hom{\mathop{\mathrm{Hom}}\nolimits}
\def\Ext{\mathop{\mathrm{Ext}}\nolimits}
\def\SheafHom{\mathop{\mathcal{H}\!\mathit{om}}\nolimits}
\def\SheafExt{\mathop{\mathcal{E}\!\mathit{xt}}\nolimits}
\def\Sch{\mathit{Sch}}
\def\Mor{\mathop{\mathrm{Mor}}\nolimits}
\def\Ob{\mathop{\mathrm{Ob}}\nolimits}
\def\Sh{\mathop{\mathit{Sh}}\nolimits}
\def\NL{\mathop{N\!L}\nolimits}
\def\CH{\mathop{\mathrm{CH}}\nolimits}
\def\proetale{{pro\text{-}\acute{e}tale}}
\def\etale{{\acute{e}tale}}
\def\QCoh{\mathit{QCoh}}
\def\Ker{\mathop{\mathrm{Ker}}}
\def\Im{\mathop{\mathrm{Im}}}
\def\Coker{\mathop{\mathrm{Coker}}}
\def\Coim{\mathop{\mathrm{Coim}}}

% Boxtimes
%
\DeclareMathSymbol{\boxtimes}{\mathbin}{AMSa}{"02}

%
% Macros for moduli stacks/spaces
%
\def\QCohstack{\mathcal{QC}\!\mathit{oh}}
\def\Cohstack{\mathcal{C}\!\mathit{oh}}
\def\Spacesstack{\mathcal{S}\!\mathit{paces}}
\def\Quotfunctor{\mathrm{Quot}}
\def\Hilbfunctor{\mathrm{Hilb}}
\def\Curvesstack{\mathcal{C}\!\mathit{urves}}
\def\Polarizedstack{\mathcal{P}\!\mathit{olarized}}
\def\Complexesstack{\mathcal{C}\!\mathit{omplexes}}
% \Pic is the operator that assigns to X its picard group, usage \Pic(X)
% \Picardstack_{X/B} denotes the Picard stack of X over B
% \Picardfunctor_{X/B} denotes the Picard functor of X over B
\def\Pic{\mathop{\mathrm{Pic}}\nolimits}
\def\Picardstack{\mathcal{P}\!\mathit{ic}}
\def\Picardfunctor{\mathrm{Pic}}
\def\Deformationcategory{\mathcal{D}\!\mathit{ef}}
\usepackage{cleveref}
\crefname{originallemma}{original lemma}{original lemmas}
\Crefname{originallemma}{Original lemma}{Original lemmas}
\makeatletter
\newlabel{lemma-f-bounded}{{0H84}{source}{}{sourcepointer-1}{}}
\newlabel{lemma-f-bounded@cref}{{[originallemma][0][]0H84}{[1][0][]source}{}{}{}}
\makeatother
\crefname{originalsection}{original section}{original sections}
\Crefname{originalsection}{Original section}{Original sections}
\makeatletter
\newlabel{section-tor}{{0651}{source}{}{sourcepointer-2}{}}
\newlabel{section-tor@cref}{{[originalsection][0][]0651}{[1][0][]source}{}{}{}}
\makeatother
\title{An object of D(R) with vanishing positive cohomology and flat localizations at f and at R over fR is a flat module in degree zero}
\author{The Stacks Project Authors}
\date{stacks-project:more-algebra.tex:local-snapshot-20261012; GNU FDL 1.2 or later}
\begin{document}
\maketitle
\setcounter{subsection}{1}
\noindent\textit{Editorial note.} All delivered mathematical text is the author's own, copied from the pinned Stacks Project chapter file more-algebra.tex; no statement or proof step was rewritten, reordered or removed. The counted statement is the labeled lemma environment at source lines 30032-30043 (canonical tag 0H85) and the counted proof is the proof environment at source lines 30045-30080, which immediately follows it and is its own proof.\par\medskip
\noindent\textit{Editorial note.} The proof refers twice inside the chapter: once to the neighbouring lemma lemma-f-bounded at source line 30012 for the flatness criterion and once to the section on Tor at source line 16853 for the meaning of flatness of a derived object. Both are navigation pointers whose roles are recorded in the external reference reasons; the two computations that discharge the hypotheses of that criterion are inside the counted span.\par\medskip
\noindent\textit{Editorial note.} The blank lines at source lines 30081-30083 separate this proof from the next statement of the chapter and are outside the counted spans.\par\medskip
\noindent\textit{Editorial note.} The chapter file pinned-original-source.tex carries only an input of the shared preamble on its first line, so no preamble text lives inside the chapter. The retained file sources/stacks-more-algebra/retained-preamble.tex is a mechanical, line-for-line copy of the project's single shared preamble as staged in the snapshot (Xy-pic setup plus the native theorem declarations and mathematical macros). Nothing was added, removed or redefined, and the excerpt serves only as preamble, never as a span coordinate.\par\medskip
\noindent\textit{Editorial note.} The source URL is the source-level Stacks Project URL, as required when more than one card is produced from one source; the per-claim tag is carried by claim\_key and by this note. The counted statement and proof cite no bibliography entry, so this card needs no bibliography.\par\medskip
\noindent\textit{Allowed standard prerequisites (editorial).}\begin{itemize}
\item A bounded above complex whose cohomology is concentrated in a single degree and vanishes modulo f and after inverting f is a flat module placed in that degree. invoked at proof lines 30049-30053, where the author reduces to the vanishing of the localizations of the cohomology and of the cohomology of the cone of multiplication by f. States the criterion in the form used: it is enough to check the two vanishing conditions for the object N = M tensor K.
\item The derived tensor product commutes with localization and with the quotient by f, so the localizations of the cohomology of N can be computed from the corresponding derived objects of M and K. used at proof lines 30054-30061 and again at proof lines 30071-30077. Computes the localizations of the cohomology of N in terms of M tensor K tensor R\_f and of the object annihilated by f, so that assumptions (2) and (3) apply.
\item For an element f of a ring the cone of multiplication by f on a complex is computed by the cone of multiplication by f on any representative, so the cone of N is N tensored with the cone of f on K. used at proof lines 30061-30070, where C = M tensor C' and the distinguished triangle with K[f] and K over fK is written. Reduces the second vanishing condition to the two-term cohomology of the cone, which is controlled by assumption (3).
\end{itemize}\medskip
\ifdefined\hypertarget\hypertarget{sourcepointer-1}{}\fi
\noindent\textit{Source reference.} Original lemma 0H84: sources/stacks-more-algebra/pinned-original-source.tex, source lines 30012--30012. The reference names the criterion that reduces flatness of N to two vanishing conditions, used at proof lines 30050-30053; the verification of both conditions for N, which is the substance of the proof, is carried out inside the counted span at lines 30054-30079.\par\medskip
\ifdefined\hypertarget\hypertarget{sourcepointer-2}{}\fi
\noindent\textit{Source reference.} Original section 0651: sources/stacks-more-algebra/pinned-original-source.tex, source lines 16853--16853. The reference points to the section in which the criterion for a derived object to be a flat module placed in degree 0 is formulated and used at proof line 30048; the paragraph of the proof containing the actual computations stays inside the counted span.\par\medskip
\begin{lemma}
\label{lemma-f-flat}
Let $R$ be a ring and $f \in R$. Let $M \in D(R)$. Assume
\begin{enumerate}
\item $H^i(M) = 0$ for $i > 0$,
\item $M \otimes_R^\mathbf{L} R_f$ is isomorphic to a flat $R_f$-module
placed in degree $0$,
\item $M \otimes_R^\mathbf{L} R/fR$ is isomorphic to a flat $R/fR$-module
placed in degree $0$.
\end{enumerate}
Then $M$ is isomorphic to a flat $R$-module placed in degree $0$.
\end{lemma}

\begin{proof}
Let $K$ be an $R$-module. It suffices to show that
$N = M \otimes_R^\mathbf{L} K$ is isomorphic to an
$R$-module placed in degree $0$, see Section \ref{section-tor}.
By assumption (1) we see that $N$ only has nonvanishing
cohomology in degrees $\leq 0$. Thus, by Lemma \ref{lemma-f-bounded},
it suffices to show that $H^i(N)_f = 0$ for $i < 0$ and
that $H^i(C) = 0$ for $i < -1$ where $C$ is the cone
of $f : N \to N$. For the first, we note that
$$
H^i(N)_f = H^i(N \otimes_R^\mathbf{L} R_f) =
H^i(M \otimes_R^\mathbf{L} K \otimes_R^\mathbf{L} R_f ) =
H^i((M \otimes_R^\mathbf{L} R_f) \otimes_{R_f}^\mathbf{L}
(K \otimes_R^\mathbf{L} R_f))
$$
By our assumption (2) this is zero, except if $i = 0$ and then one gets
$H^0(M)_f \otimes_R K$. For the second, let $C'$ be the cone of $f : K \to K$.
Then
$$
C = M \otimes_R^\mathbf{L} C'
$$
Now $C'$ only has nonzero cohomology in degrees $0$ and $-1$ equal to
$K/fK$ and $K[f]$; in other words, there is a distinguished triangle
$$
(K[f])[1] \to C' \to K/fK
$$
On the other hand, for any $R$-module $K'$ annihilated by $f$ we have
$$
M \otimes_R^\mathbf{L} K' =
M \otimes_R^\mathbf{L} R/fR \otimes_{R/fR}^\mathbf{L} K'
$$
By our assumption (3) this is equal to the module
$H^0(M \otimes_R^\mathbf{L} R/fR) \otimes_{R/fR} K'$ placed in degree $0$.
Combining the above we conclude that $C$ only has nonzero
cohomology in degrees $0$ and $-1$ and the proof is complete.
\end{proof}

\end{document}
